\subsection*{v11 --- review-pass closure and typed targets for
the remaining upstream theorems (the present revision)}

Relative to v10, this version is a closure-and-naming pass. It
addresses the v10 reviewer's three actionable findings:

\begin{itemize}[leftmargin=1.6em, itemsep=2pt]
\item \textbf{Editorial consistency.}\ The stale "Stage D
    characterisation as targeted future work" paragraph and the
    matching "Stage D backward" entry in \S\ref{sec:next}, both
    inherited from v9, are removed: v10's
    \texttt{YukawaStageDExistence} closed that gap.

\item \textbf{Wording.}\ "Upstream interfaces structurally
    closed" is reworded to "closed to the next algebraic
    layer", honestly reflecting that Calder\'on, Higgs Stage B,
    and Yukawa Stage D backward are now Lean-proved
    one-layer-up structural derivations, while the analytic
    / TFPT-internal content of \emph{that next layer up}
    remains open.

\item \textbf{Three-part claim contract.}\ The
    \texttt{tfptclaimcontract} environment is restructured to
    separate (a) direct \texttt{CarrierPremises} inputs,
    (b) Lean-proved upstream refinements, (c) open
    TFPT-origin theorems. The third group now lists three
    typed Lean targets for the open content.

\item \textbf{New module
    \texttt{BoundaryPolarization.ofCarrierInvolution}.}\
    The converse of v10's
    \texttt{BoundaryPolarization.toCarrierInvolution}: given
    a carrier involution $\varepsilon^2 = \mathbf 1$ on
    \texttt{Module.End K H}, the +1 and -1 eigenspaces form a
    \texttt{BoundaryPolarization}. This makes the Paper-1
    interface a Lean-proved \emph{algebraic equivalence}: the
    polarisation picture and the involution picture are two
    presentations of the same finite-dimensional carrier
    datum. \emph{Not} the analytic Paper 1 derivation; the
    construction of the polarisation from boundary elliptic
    theory remains upstream.

\item \textbf{New module \texttt{SeamWindingInterface.lean}.}\
    Introduces \texttt{SeamWindingData}\,$(m, n)$, an abstract
    typed target for TFPT Paper 2's seam-class theorem
    "$[u_\Sigma] = 1$ produces a primitive oriented
    determinant-preserving carrier". Provides bidirectional
    bridges to and from \texttt{PrimitiveOrientedDeterminantCarrier}:
    the seam-winding interface is the existing rigidity input,
    renamed to make the upstream TFPT theorem visible.

\item \textbf{New module \texttt{BoundaryYukawaKernelInterface.lean}.}\
    Introduces \texttt{BoundaryYukawaKernel}\,$K\,E$,
    \emph{packaging}\ the Yukawa Stage D quadruple
    $(\omega, \texttt{CI}, \texttt{CS}, \texttt{Nontrivial})$.
    This is \emph{not} a new physical derivation; it is a
    typed renaming of the Stage D condition, supplied so that
    the open TFPT theorem
    "boundary kernel produces a canonical $\omega$" has an
    explicit Lean-side conclusion. Forward direction
    \texttt{finrank\_E\_eq\_three} reuses v10's Stage D
    discharge; backward direction \texttt{ofFinrankEqThree}
    reuses v10's \texttt{YukawaStageDExistence} to give the
    characterisation
    \texttt{nonempty\_iff\_finrank\_eq\_three}. The
    \emph{canonicity} of $\omega$ from the boundary kernel
    --- whether the TFPT-internal construction produces the
    same form as the basis-built witness, up to carrier
    symmetries --- remains the upstream physics question.

\item \textbf{AuditCheck, AxiomCheck, AuditContract extended.}\
    Fifteen new headline names added across the three new
    modules; all depend only on
    $\{\texttt{propext},\ \texttt{Classical.choice},\
    \texttt{Quot.sound}\}$. \texttt{audit.sh} returns
    \texttt{AUDIT:\ PASS}.
\end{itemize}

\smallskip
After v11, the boundaries of the Lean-side carrier-rigidity
chain are sharply visible. The three open TFPT-origin theorems
each have a typed Lean target:

\[
\begin{aligned}
\text{Paper 1 analytic boundary theory}
&\;\Rightarrow\; \texttt{BoundaryPolarization}, \\
\text{TFPT boundary kernel + reps}
&\;\Rightarrow\; \texttt{BoundaryYukawaKernel}, \\
\text{seam class }[u_\Sigma] = 1
&\;\Rightarrow\; \texttt{SeamWindingData}.
\end{aligned}
\]

\subsection*{v10 --- upstream interfaces closed to the next
algebraic layer}

Relative to v9, this version closes the three standing upstream
gaps documented in the v9 limitations section. Each gap is
replaced by a Mathlib-native structural Lean module:

\begin{itemize}[leftmargin=1.6em, itemsep=2pt]
\item \textbf{New module \texttt{BoundaryPolarization.lean}.}\
    Introduces the structure
    \texttt{BoundaryPolarization K H := \{Hplus Hminus :
    Submodule K H, isCompl Hplus Hminus\}}, the algebraic
    skeleton of the Paper 1 polarisation. The projection
    onto $H_+$ along $H_-$ is an idempotent endomorphism
    (via Mathlib's \texttt{Submodule.IsCompl.projection} and
    its \texttt{IsIdempotentElem} lemma); a bridge
    \texttt{toCalderonProjector} feeds the
    \texttt{CalderonProjector} chain in
    \texttt{Module.End K H}, using the standard
    \texttt{Invertible (2 : \texttt{End} K M)} instance from
    Mathlib's \texttt{QuadraticForm.Basic}.

\item \textbf{New module \texttt{HiggsSchemeCohomologyShadow.lean}.}\
    Constructs a \texttt{LinearEquiv} between
    \texttt{DegreeOneTwoVarForm K} (the algebraic shadow used
    by \texttt{HiggsTopForm}) and Mathlib's
    \texttt{MvPolynomial.homogeneousSubmodule (Fin 2) K 1}, the
    degree-1 piece of the bivariate polynomial ring. The latter
    is the algebraic-geometry "Proj-side" presentation of
    $H^0(\mathbb P^1, \mathcal O(1))$. The proof uses
    \texttt{MvPolynomial.homogeneousSubmodule\_one\_eq\_span\_X}
    and \texttt{Submodule.mem\_span\_range\_iff\_exists\_fun}
    to extract coefficients; injectivity is by direct
    extraction of $X_i$-coefficients.

\item \textbf{New module \texttt{YukawaStageDExistence.lean}.}\
    Closes the Stage D backward direction:
    $\dim E = 3 \Rightarrow \exists\,\omega$ with
    \texttt{ContractionInjective} and
    \texttt{ContractionSurjective}. The construction uses
    Mathlib's \texttt{Module.Basis.det} (the determinant of
    vectors in a basis) transported via
    \texttt{exteriorPower.alternatingMapLinearEquiv} into a
    Yukawa form on $\Lambda^3 E$. Injectivity is proved by
    direct $3\times 3$ matrix-determinant computation in
    the basis; surjectivity follows from
    $\dim E = \dim(\Lambda^2 E)^* = 3$ and Mathlib's
    \texttt{LinearMap.injective\_iff\_surjective\_of\_finrank\_eq\_finrank}.
    The Stage D characterisation
    $\exists\,\omega\;(\textrm{CI}\,\wedge\,\textrm{CS})
      \iff \dim E = 3$ now holds (under
    \texttt{Nontrivial}\,$E$).

\item \textbf{AuditCheck, AxiomCheck, AuditContract extended.}\
    Eighteen new headline names added; all depend only on
    $\{\texttt{propext},\ \texttt{Classical.choice},\
    \texttt{Quot.sound}\}$. \texttt{audit.sh} returns
    \texttt{AUDIT:\ PASS} with seven hard checks.
\end{itemize}

\smallskip
After v10, the conditional Lean theorem is

\[
\boxed{\;
\begin{aligned}
\texttt{BoundaryPolarization}\;(H = H_+ \oplus H_-)
&\;\Rightarrow\; \pi^2 = \pi
\;\Rightarrow\; \varepsilon^2 = \mathbf 1
\;\Rightarrow\; P_\pm \text{ idempotent} \\
K[X_0, X_1]_1 \simeq \texttt{DegreeOneTwoVarForm}
&\;\Rightarrow\; \dim E_+ = 2 \\
\dim E_- = 3\;(\text{with }\texttt{Nontrivial}\,E_-)
&\;\iff\; \exists\,\omega\;(\textrm{CI} \wedge \textrm{CS}) \\
m q_- + n q_+ = 0,\ q_- < 0 < q_+,\ \gcd = 1
&\;\Rightarrow\; (q_-, q_+) = (-n/d, m/d)
\end{aligned}
\;}
\]

with every step a Lean theorem, no \texttt{sorry}, all axioms
$\subseteq\{\texttt{propext},\ \texttt{Classical.choice},\
\texttt{Quot.sound}\}$.

\subsection*{v9 --- review-readiness pass and Stage D transparency}

Relative to v8, this version is a consistency-and-transparency
revision rather than a structural one. It addresses internal
version drift identified in the v8 review and adds one
unconditional transparency theorem at the Stage D interface:

\begin{itemize}[leftmargin=1.6em, itemsep=2pt]
\item \textbf{New theorem
    \texttt{YukawaTrilinearForm.finrank\_E\_eq\_zero\_or\_three}.}\
    For any finite-dimensional $K$-vector space (no
    \texttt{Nontrivial} hypothesis) a trilinear form with
    \texttt{ContractionInjective} and
    \texttt{ContractionSurjective} satisfies
    $\dim E = 0 \lor \dim E = 3$. The two predicates are
    vacuously satisfied at $E = 0$, which is the mathematical
    reason the existing rank-3 theorem requires
    \texttt{Nontrivial}\,$E$.

\item \textbf{Stage D slogans and statements carry
    \texttt{Nontrivial}\,$E$ explicitly.}\ The frontbox,
    abstract, headline boxed equations, and interface table now
    state the rank-3 conclusion under \texttt{Nontrivial}\,$E$.
    The Lean theorem already carried this hypothesis; v9 fixes
    the documentation to match.

\item \textbf{Predicate naming unified throughout the note.}\
    All operational mentions use \texttt{ContractionInjective}
    and \texttt{ContractionSurjective}. The TFPT-side phrases
    "non-degenerate" and "primitive / no spectator" are
    retained only as docstring synonyms in mathematical prose.

\item \textbf{Audit gate uniformly described as seven checks.}\
    The earlier v6 wording "six checks" is removed; the
    \texttt{AuditContract.lean} signature-lock check is
    consistently the seventh.

\item \textbf{Limitations section rewritten to match v8/v9.}\
    The previous "Higgs Stage A only" and "Calder\'on interface
    is documentary only" wording is replaced by the accurate v9
    status: Higgs has a Stage B algebraic shadow, Calder\'on has
    an algebraic idempotent-to-involution layer, and the
    remaining open content is precisely identified.

\item \textbf{Calder\'on and Higgs wording sharpened.}\ The note
    no longer claims that Paper 1 or scheme cohomology is
    formalised. Calder\'on is presented as the algebraic step
    $\pi^2 = \pi \Rightarrow \varepsilon^2 = \mathbf 1$; Higgs is
    presented as the algebraic shadow $E_+ \simeq
    \texttt{DegreeOneTwoVarForm K}$.

\item \textbf{AxiomCheck, AuditCheck, AuditContract updated.}\
    The three audit modules now include the new
    \texttt{finrank\_E\_eq\_zero\_or\_three} headline. All
    audit checks continue to return \texttt{AUDIT:\ PASS}.
\end{itemize}

\subsection*{v8 --- Calderón structural + Higgs Stage B}

Relative to v7, this version upgrades two upstream interfaces
from documentary wrappers to structural derivations, narrowing
the "out of scope" column of the interface table:

\begin{itemize}[leftmargin=1.6em, itemsep=2pt]
\item \textbf{New module \texttt{CalderonProjector.lean}.}\
    Introduces the structure
    \texttt{CalderonProjector A := \{\(\pi : A\), \(\pi^2 = \pi\)\}},
    the algebraic skeleton of the Paper 1 Calder\'on
    construction. Defines the spectral reflection
    $\varepsilon := 2\pi - 1$ and \emph{proves}
    $\varepsilon^2 = \mathbf 1$ via the elementary identity
    $(2\pi - 1)^2 = 4\pi^2 - 4\pi + 1 = 1$. The bridge
    \texttt{toCalderonCertificate} feeds the structural result
    into the existing
    \texttt{CalderonInterface.CalderonCertificate} wrapper, so
    the Paper 1 input is now derived from a single algebraic
    hypothesis: \emph{the existence of an idempotent}.

\item \textbf{New module \texttt{HiggsTopForm.lean}.}\
    Introduces the structure
    \texttt{HiggsTopForm K Eplus :=
            \{Eplus $\simeq$ DegreeOneTwoVarForm K\}},
    where
    \texttt{DegreeOneTwoVarForm K := \{coeffX coeffY : K\}}
    is the algebraic shadow of
    $H^0(\mathbb P^1_K, \mathcal O(1))$ --- the space of
    homogeneous degree-1 polynomials $\alpha X + \beta Y$ in
    two formal variables over $K$. The Stage-B rank theorem
    $\dim E_+ = 2$ follows mechanically (\texttt{finrank\_eq\_two}
    on \texttt{DegreeOneTwoVarForm}, transported via the
    \texttt{HiggsTopForm} equivalence). A bridge
    \texttt{toHiggsIndexCertificate} produces the Stage A
    witness.

\item \textbf{Updated audit gate.}\
    \texttt{AxiomCheck.lean}, \texttt{AuditCheck.lean}, and
    \texttt{AuditContract.lean} extended with five new headline
    theorems. All five depend only on
    $\{\texttt{propext},\ \texttt{Classical.choice},\
    \texttt{Quot.sound}\}$. \texttt{audit.sh} continues to
    return \texttt{AUDIT:\ PASS} with seven hard checks.
\end{itemize}

\smallskip
The Lean-side picture from v8 onwards is:

\[
\begin{aligned}
\text{Calder\'on projector}\; (\pi^2 = \pi)
&\;\Rightarrow\; \varepsilon^2 = \mathbf 1 \;\Rightarrow\; P_-, P_+ \text{ idempotent}\\
\text{Higgs Stage B}\; (E_+ \simeq \alpha X + \beta Y)
&\;\Rightarrow\; \dim E_+ = 2\\
\text{Yukawa Stage D}\; (\omega, \texttt{CI}, \texttt{CS}, \texttt{Nontrivial}\,E_-)
&\;\Rightarrow\; \dim E_- = 3
\end{aligned}
\]

with all three discharges now derived from \emph{structural
hypotheses}, not bare numerical / linear-equivalence
assumptions.

\paragraph{Remaining Lean lemma gap.}
The Stage D characterisation backward direction --- existence
of a $(\omega, \texttt{ContractionInjective},
\texttt{ContractionSurjective})$ triple from $\dim E = 3$ ---
requires showing that the wedge map
$E \to \mathrm{Hom}(\Lambda^2 E, \Lambda^3 E)$, $v \mapsto (\eta \mapsto v \wedge \eta)$,
is injective in dimension 3. The cleanest formalisation route
extends a non-zero vector to a basis and uses
\texttt{exteriorPower.iotaMulti\_family\_linearIndependent\_ofBasis}
to deduce $v \wedge e_1 \wedge e_2 \neq 0$. Mathlib does not
yet have a direct lemma for the non-degeneracy of the wedge
pairing $\Lambda^p E \otimes \Lambda^{n-p} E \to \Lambda^n E$,
so this is a single missing infrastructure lemma; the
remainder of the proof is bookkeeping.

\subsection*{v7 --- Signature lock, predicate rename, interface
discipline}

Relative to v6, this version hardens the audit interface,
neutralises a semantically-loaded naming choice, and tightens
the editorial framing around what the note does and does not
claim:

\begin{itemize}[leftmargin=1.6em, itemsep=2pt]
\item \textbf{New \texttt{AuditContract.lean} module.}\ Each
    headline theorem is now used in an
    \texttt{example : <exact intended type> := <theorem>}
    pattern. This is strictly stronger than the \texttt{\#check}
    audit of \texttt{AuditCheck.lean}: it detects silent
    weakening of theorem statements that \texttt{\#check} would
    miss (since a weakened theorem still elaborates under its
    own name). The \texttt{audit.sh} gate now has \emph{seven}
    checks; the new seventh requires
    \texttt{AuditContract.lean} to elaborate.

\item \textbf{Stage D predicate rename.}\
    \texttt{YukawaTrilinearForm.Nondegenerate} $\to$
    \texttt{ContractionInjective} and
    \texttt{YukawaTrilinearForm.Primitive} $\to$
    \texttt{ContractionSurjective}. The new names are purely
    operational --- they describe the bijectivity of
    $C_\omega$ --- and leave the semantic content
    ("non-degenerate", "primitive") to the docstrings. A
    combined predicate
    \texttt{PrimitiveIndecomposableYukawaCondition} is
    available for the TFPT-side phrase. This removes the
    "primitive renamed to surjective" attack surface.

\item \textbf{New
    \texttt{OrientedDeterminantCarrier.lean} module.}\
    The three-part TFPT carrier condition --- trace-zero,
    orientation, primitivity --- is now packaged as a single
    Lean structure
    \texttt{PrimitiveOrientedDeterminantCarrier (m n : \(\mathbb N\))}.
    Producing such a structure is the formal target for
    a future TFPT-internal theorem
    "seam winding $[u_\Sigma] = 1$ implies a primitive oriented
    determinant-preserving carrier".

\item \textbf{Interface table.}\ A new
    \S\ref{sec:interface-table} maps every load-bearing TFPT
    Paper 2 claim to its current Lean status and the missing
    upstream proof. This separates "formalised", "interfaced",
    and "out of scope" without ambiguity.

\item \textbf{Claim Contract refinement.}\ The contract now
    distinguishes \emph{direct} inputs to
    \texttt{CarrierPremises} (the rank witnesses and the
    boundary involution) from \emph{refinable} upstream
    interfaces (Stages B/C/D feeding the Yukawa rank witness).

\item \textbf{Editorial fixes.}\ The running header is
    overridden for this note (it previously inherited the
    TFPT 4.5 series header). The "v3" relic in the Result
    section is corrected.
\end{itemize}

